\documentclass[UTF8,openany,11pt,a4paper]{ctexbook}
\usepackage{fontspec}
\setCJKmainfont{Noto Serif CJK SC}
\setCJKsansfont{Noto Sans CJK SC}
\setmainfont{Latin Modern Roman}
\usepackage[margin=24mm,headheight=15pt]{geometry}
\usepackage{amsmath,amssymb,mathtools,mhchem,booktabs,longtable}
\usepackage{microtype,xcolor,tcolorbox,fancyhdr,hyperref}
\definecolor{ink}{HTML}{193247}
\definecolor{accent}{HTML}{087F76}
\definecolor{soft}{HTML}{F0F7F6}
\definecolor{warn}{HTML}{8A4513}
\hypersetup{colorlinks=true,linkcolor=ink,urlcolor=accent,pdfauthor={Ruoshui Chemistry Editorial Review},pdftitle={大学四大化学·原创高阶综合题科学复核册}}
\setlength{\parindent}{2em}
\setlength{\parskip}{.25em}
\linespread{1.16}
\setcounter{secnumdepth}{2}
\pagestyle{fancy}
\fancyhf{}
\fancyhead[L]{\small 三千弱水 \quad 化学综合题科学复核}
\fancyhead[R]{\small \textcolor{warn}{未通过专家审定}}
\fancyfoot[C]{\thepage}
\renewcommand{\headrulewidth}{.35pt}
\newtcolorbox{reviewnote}{colback=soft,colframe=accent!60!black,arc=1.3mm,boxrule=.55pt,left=2mm,right=2mm}
\newcommand{\sciencewarn}[1]{\begin{reviewnote}\textbf{\textcolor{warn}{模型适用条件与审核重点：}}#1\end{reviewnote}}
\begin{document}
\begin{titlepage}
\vspace*{22mm}
{\color{accent}\rule{\linewidth}{2.3pt}}\\[12mm]
{\Large 三千弱水 · 大学四大化学题库}\\[12mm]
{\Huge\bfseries 原创高阶综合题}\\[4mm]
{\Huge\bfseries 科学复核册}\\[12mm]
{\large 第一辑 \quad 12 道独立专题 \quad 题目与完整推导}\\[16mm]
\begin{reviewnote}
本册供\textbf{科学性、题意条件、计算和排版}联合复核。
全部题目为原创教学题，不宣称出自真实大学考试；
无外部专家签名，不作为已审定的正式试题册发布。
特别区分机器方程残差为零与化学模型真实可靠两种概念。
\end{reviewnote}
\vfill
{\small 版本：2026-10-28 \quad 校验：数值反算、来源许可、假设清单}\\[6mm]
{\small 许可证：独立原创题文 CC0-1.0；第三方材料不在本声明范围内}
{\color{accent}\rule{\linewidth}{1.1pt}}
\end{titlepage}
\frontmatter
\tableofcontents
\mainmatter
\chapter*{编者说明}
\addcontentsline{toc}{chapter}{编者说明}
题库现有大量用于计算训练的参数变式，本辑不继续通过数字替换扩大名义独立知识点。
本册选入\textbf{每科三道}独立撰写的高阶综合题，强调建模、约束条件、数值求解、
守恒检验、近似适用范围以及易错点。个别题目使用明确标记的理想化模型或假想体系；
这不是现实实验结果。机器脚本可以重算表内结果，仍需人工科学审阅。
\part{题目册}
\chapter{无机化学}
\section{I01 \quad 闭合氨配位体系中 AgCl 的非线性溶解}
在 \(1.000\,\mathrm L\) 的理想溶液中有过量固体 \(\ce{AgCl(s)}\)，体系的游离氨及络合氨总分析浓度为 \(C_{\mathrm L}=0.1000\,\mathrm{mol\,L^{-1}}\)。仅允许 \(\ce{Ag+}\)、\(\ce{[Ag(NH3)2]+}\)、\(\ce{Cl-}\) 及游离 \(\ce{NH3}\) 存在，并假设 \(\ce{NH3}/\ce{NH4+}\) 酸碱平衡及其它配合物可忽略。给定 \(K_{\rm sp}=1.80\times10^{-10}\) 和 \(\beta_2=1.60\times10^7\,\mathrm{L^2\,mol^{-2}}\)，同时求游离氨、游离银、二氨合银和总溶解银浓度；不得把游离氨浓度直接等同于 \(C_{\mathrm L}\)。
\sciencewarn{理想活度；明确忽略 NH4+、Ag(NH3)+、其他络合及溶解产生的体积改变；过量 AgCl 固体仍存在。}
\section{I02 \quad 有限体积原电池的 Nernst 电势与平衡转化}
假想二价金属电池反应为 \(\ce{M(s) + N^{2+}(aq) <=> M^{2+}(aq) + N(s)}\)，在 \(298.15\,\mathrm K\) 时标准电动势 \(E^\circ=0.0450\,\mathrm V\)。同一 \(0.2500\,\mathrm L\) 理想、均匀的反应体系初始 \([\ce{M^{2+}}]_0=0.0100\,\mathrm M\)、\([\ce{N^{2+}}]_0=0.0200\,\mathrm M\)，金属固体过量。忽略体积改变和液接电势。求初始电动势、无量纲平衡常数、平衡反应进度（以 \(\mathrm{mol\,L^{-1}}\) 表示）及从初态到平衡传递的总电量。
\sciencewarn{假想可逆金属体系；固体活度为 1；离子活度按理想标准化浓度；电势随反应进行变化；足量固体。}
\section{I03 \quad X 射线晶胞测定、密度与空位缺陷联合反演}
某银晶体可近似为 FCC 晶格，使用波长 \(\lambda=0.1540\,\mathrm{nm}\) 的 X 射线观察到一阶 \((111)\) 衍射，给定 Bragg 角 \(\theta=18.90^\circ\)（不是 \(2\theta\)）。银摩尔质量 \(M=107.8682\,\mathrm{g\,mol^{-1}}\)，实测体密度 \(10.20\,\mathrm{g\,cm^{-3}}\)。假设晶胞边长不受缺陷影响，唯一缺陷是均匀、随机的银原子位点空位，求晶面间距、晶胞边长、满占据理想密度、空位位点分数。
\sciencewarn{晶体为理想 FCC；一阶 (111) 峰；角度是 θ；仅空位降低平均晶胞原子数，忽略夹杂、孔隙、仪器误差与应变。}
\chapter{分析化学}
\section{A01 \quad 双波长光谱解混、稀释反算与试样质量}
两种仅含 \(X,Y\) 的分析物在光程 \(1.00\,\mathrm{cm}\) 下服从 Beer--Lambert 定律，校正空白后的两波长吸光度分别为 \(A_1=0.4600\)、\(A_2=0.3400\)。摩尔吸光系数矩阵为
\[
\begin{pmatrix}10000&2000\\3000&8000\end{pmatrix}
\ \mathrm{L\,mol^{-1}cm^{-1}}.
\]
测量液由 \(10.00\,\mathrm{mL}\) 原样定容至 \(100.0\,\mathrm{mL}\) 得到。求测量液及原样的两组分浓度；再求原样 \(25.00\,\mathrm{mL}\) 中 \(X,Y\) 的质量（取 \(M_X=150.0\)、\(M_Y=200.0\,\mathrm{g\,mol^{-1}}\)）。
\sciencewarn{两种发色团彼此不反应、光程一致、吸光系数已在测量条件下标定；完全回收与线性吸收；两个测量波长可逆矩阵。}
\section{A02 \quad 双指示剂滴定反演碳酸盐组成与离子强度}
从一份仅由 \(\ce{Na2CO3}\) 与 \(\ce{NaHCO3}\) 组成的水溶液中取 \(25.00\,\mathrm{mL}\)，用 \(0.1000\,\mathrm M\) 的 \(\ce{HCl}\) 滴定。到酚酞终点累计消耗 \(8.20\,\mathrm{mL}\)，到甲基橙终点累计消耗 \(21.60\,\mathrm{mL}\)。认为两终点理想精确、无 \(\ce{OH-}\) 或其它缓冲组分。求原液两种盐的形式浓度、原液 \(250.0\,\mathrm{mL}\) 中两种盐的质量，以及按完全解离形式计算的离子强度。
\sciencewarn{以形式组分为定义的浓度；终点不计碳酸溶解/逸散损失；电中性及完全解离，仅计算形式离子强度。}
\section{A03 \quad 钙镁络合差量滴定、空白校正与总硬度}
将水样平行取 \(50.00\,\mathrm{mL}\) 两份，以 \(0.01000\,\mathrm M\) EDTA 滴定。钙专测原始读数 \(17.60\,\mathrm{mL}\)，钙镁总测原始读数 \(28.40\,\mathrm{mL}\)，两测定各有 \(0.20\,\mathrm{mL}\) 空白消耗。取 \(M_{\rm Ca}=40.078\)、\(M_{\rm Mg}=24.305\)、\(M_{\ce{CaCO3}}=100.09\,\mathrm{g\,mol^{-1}}\)。假设没有竞争金属、有效络合计量均为 \(1:1\)，求两种金属的浓度、若将整个 \(0.5000\,\mathrm L\) 同浓度样品完全回收所得各金属质量、以 \(\ce{CaCO3}\) 计的总硬度。
\sciencewarn{平行等体积取样；各自单独扣空白；Ca 专测时 Mg 完全屏蔽；总滴定无其它金属干扰；络合终点与回收率理想。}
\chapter{有机化学}
\section{O01 \quad 连续一级反应的中间体最佳分离时刻}
在封闭、恒温、恒容条件下，假想有机反应网络为 \(\ce{A -> B -> C}\)，两步严格遵守一级动力学，速率常数 \(k_1=0.1000\,\mathrm{h^{-1}}\)、\(k_2=0.0400\,\mathrm{h^{-1}}\)。初始只有 \(A\)，总物质的量 \(0.2500\,\mathrm{mol}\)，假设 \(1:1:1\) 化学计量、无逆反应和其它副反应，\(M_B=150.0\,\mathrm{g\,mol^{-1}}\)。求 \(B\) 的浓度（或摩尔分数）函数、使 \(B\) 最大的时间、该时刻三种物质的摩尔分数与 \(B\) 的理论质量。
\sciencewarn{严格的一级不可逆连续反应；正确定义 k1,k2；所有物种 1:1 转化，恒温恒容，极值时瞬时完美分离。}
\section{O02 \quad 非等当量缩聚的 Carothers 关系、脱水守恒与数均分子量}
理想双官能 AA/BB 单体缩聚体系中，初始 A 类官能团为 \(0.9600\,\mathrm{mol}\)，B 类官能团为 \(1.000\,\mathrm{mol}\)。A 类官能团反应分数为 \(p=0.9700\)；每形成一条连接键均损失一分子水。假设体系初始分子全为双官能单体、无环化、无交联、无单官能杂质；全部单体按数目平均的摩尔质量为 \(130.0\,\mathrm{g\,mol^{-1}}\)，\(M_{\ce{H2O}}=18.015\,\mathrm{g\,mol^{-1}}\)。求初始单体分子物质的量、缩合键物质的量、最终链物质的量、数均聚合度、产物总质量和数均摩尔质量。
\sciencewarn{AA/BB 双官能、无环化、无交联、每反应键脱去一分子水；p 是少量 A 官能团的转化分数；单体数均质量与起始总数相配。}
\section{O03 \quad 双响应因子手性色谱的 ee、质量平衡与收率}
某反应仅生成互为对映体的 \(R,S\) 两种产物。手性色谱中，经独立质量标准校准的检测响应满足 \(A_R=f_Rm_R\)、\(A_S=f_Sm_S\)，其中 \(A_R=1500\)、\(A_S=700\)，响应因子 \(f_R=1.20\)、\(f_S=0.800\ \mathrm{area\,mg^{-1}}\)。两对映体摩尔质量相同，没有其它产物或杂质。若理论完全转化可以得到产物总质量 \(2500\,\mathrm{mg}\)，求两对映体质量、主对映体质量百分数、正的 ee 和总分离收率。
\sciencewarn{独立校准响应因子按质量标定、两组分完整分离、对映体摩尔质量相同、无第三组分、理论产物量可信。}
\chapter{物理化学}
\section{P01 \quad 双层扩散膜的等效阻力、通量与累积转移质量}
在温度固定、稳态一维扩散中，两储液侧浓度恒定为 \(c_{\mathrm L}=2.00\)、\(c_{\mathrm R}=0.200\,\mathrm{mol\,m^{-3}}\)，中间串联两层膜：\(L_1=80.0\,\mathrm{\mu m},D_1=1.00\times10^{-10}\,\mathrm{m^2s^{-1}},K_1=0.500\)；\(L_2=120.0\,\mathrm{\mu m},D_2=3.00\times10^{-10}\,\mathrm{m^2s^{-1}},K_2=0.800\)。两侧与膜界面的局部分配遵守 \(c_{\mathrm{membrane},i}=K_i c_{\mathrm{aqueous}}\)，界面无传质阻力。膜面积 \(2.00\,\mathrm{cm^2}\)，维持上述储液浓度 \(2.00\,\mathrm h\)。求等效传质阻力、稳态物质通量、通过膜的物质的量及若 \(M=180.0\,\mathrm{g\,mol^{-1}}\) 时的质量（mg）。
\sciencewarn{稳态 1D 串联扩散；分配系数定义一致；界面局部热力学平衡；储液浓度恒定；无对流、反应与膜内蓄积。}
\section{P02 \quad 热容差修正反应热力学与平衡常数}
某假想标准反应在 \(T_0=298.15\,\mathrm K\) 的标准焓、标准熵分别为 \(\Delta_rH^\circ(T_0)=45.000\,\mathrm{kJ\,mol^{-1}}\) 与 \(\Delta_rS^\circ(T_0)=120.00\,\mathrm{J\,mol^{-1}K^{-1}}\)。在 \(298.15\) 到 \(350.00\,\mathrm K\) 区间，标准摩尔热容差恒为 \(\Delta_rC_p^\circ=50.00\,\mathrm{J\,mol^{-1}K^{-1}}\)。求 \(350.00\,\mathrm K\) 的标准焓、熵、Gibbs 能和无量纲热力学平衡常数 \(K\)，并给出该温度的 van't Hoff 导数 \(d\ln K/dT\)。
\sciencewarn{ΔCp 在区间内常数；T 为绝对温度；K 由统一标准态活度定义并无量纲；近似只用于指定温度区间。}
\section{P03 \quad 有限溶液 Langmuir 吸附的非线性质量平衡}
把 \(m=2.000\,\mathrm g\) 的均匀吸附剂加入 \(V=0.5000\,\mathrm L\)、初浓度 \(C_0=0.002000\,\mathrm{mol\,L^{-1}}\) 的单溶质溶液中。吸附平衡遵守 \(q(C)=q_{\max}KC/(1+KC)\)，其中 \(q_{\max}=4.00\times10^{-4}\,\mathrm{mol\,g^{-1}}\)、\(K=500.0\,\mathrm{L\,mol^{-1}}\)。假设没有化学反应或容器吸附且溶液体积不变，求平衡游离浓度 \(C\)、表面覆盖度 \(\theta\)、吸附总物质的量及去除率；必须考虑吸附导致的溶液浓度下降。
\sciencewarn{单一吸附溶质；温度不变；理想单层 Langmuir、独立同质位点；没有溶液体积变化、容器损失和其它反应。}
\cleardoublepage\part{严格解答与计算验证}
\chapter{无机化学：解答}
\section{I01 \quad 闭合氨配位体系中 AgCl 的非线性溶解}
令 \(L=[\ce{NH3}],\ p=[\ce{Ag+}],\ z=[\ce{[Ag(NH3)2]+}],\ s=p+z=[\ce{Cl-}]\)。电中性由此自动满足。沉淀平衡、配合平衡及氨物料守恒为
\[
ps=K_{\rm sp},\qquad z=\beta_2pL^2,\qquad L+2z=C_{\mathrm L}.
\]
消去 \(p,z,s\)，得
\[
p=\sqrt{\frac{K_{\rm sp}}{1+\beta_2L^2}},\quad
z=\beta_2L^2\sqrt{\frac{K_{\rm sp}}{1+\beta_2L^2}},
\quad
F(L)=L+2z-C_{\mathrm L}=0.
\]
由于 \(z(L)=\sqrt{K_{\rm sp}}\,\beta_2L^2(1+\beta_2L^2)^{-1/2}\) 对 \(L>0\) 严格递增，且 \(F(0)<0<F(C_{\mathrm L})\)，因此物理解唯一。用二分法求得
\[
\begin{aligned}
L&=0.09030725\,\mathrm M, &
p&=3.714086\times10^{-8}\,\mathrm M,\\
z&=0.004846377\,\mathrm M, &
s&=0.004846414\,\mathrm M.
\end{aligned}
\]
将数值代回可得 \(ps=1.80\times10^{-10}\)、\(z/(pL^2)=1.60\times10^7\)、\(L+2z=0.1000\,\mathrm M\)。
\begin{reviewnote}
本题计算机独立数值见证：1：0.090307246；2：3.7140864e-8；3：0.0048463768；4：0.0048464140。
逆向方程最大相对残差：2.776e-16。
\end{reviewnote}
\section{I02 \quad 有限体积原电池的 Nernst 电势与平衡转化}
反应商 \(Q=[\ce{M^{2+}}]/[\ce{N^{2+}}]\) 的数值采用相同标准浓度归一化。Nernst 方程为
\[
E=E^\circ-\frac{RT}{2F}\ln Q,\qquad
K=\exp\!\left(\frac{2FE^\circ}{RT}\right).
\]
初始 \(Q_0=0.0100/0.0200=0.5\)，故 \(E_0=0.0539044\,\mathrm V\)。在平衡时令反应进度按单位体积记为 \(\xi\)，则
\[
\frac{0.0100+\xi}{0.0200-\xi}=K,\quad
\xi=\frac{0.0200K-0.0100}{K+1}.
\]
代入 \(R=8.314462618\,\mathrm{J\,mol^{-1}K^{-1}}\)、\(F=96485.33212\,\mathrm{C\,mol^{-1}}\)，得
\[
K=33.2135,\qquad \xi=0.01912315\,\mathrm M.
\]
因此平衡 \([\ce{M^{2+}}]=0.02912315\,\mathrm M\)、\([\ce{N^{2+}}]=0.000876846\,\mathrm M\)。转移电荷为
\[
Q_{\rm el}=2F\cdot \xi \cdot (0.2500\,\mathrm L)=922.552\,\mathrm C.
\]
注意电荷计算中 \(\xi\) 是浓度形式，必须乘溶液体积。
\begin{reviewnote}
本题计算机独立数值见证：1：0.053904369；2：33.213519；3：0.019123154；4：922.55192。
逆向方程最大相对残差：2.781e-15。
\end{reviewnote}
\section{I03 \quad X 射线晶胞测定、密度与空位缺陷联合反演}
由一阶 Bragg 定律和立方晶系面间距关系，
\[
2d_{111}\sin\theta=\lambda,\qquad
d_{111}=\frac{a_{\rm cell}}{\sqrt3}.
\]
故
\[
d_{111}=0.2377149\,\mathrm{nm},\qquad
a_{\rm cell}=0.4117343\,\mathrm{nm}.
\]
FCC 常规晶胞有 \(4\) 个可能原子位点，若全部被占据，其密度为
\[
\rho_0=\frac{4M}{N_Aa_{\rm cell}^{\,3}}
       =10.26482\,\mathrm{g\,cm^{-3}},
\]
其中 \(a_{\rm cell}(\mathrm{cm})=a_{\rm cell}(\mathrm{nm})\times10^{-7}\)。若平均空位分数为 \(v\)，有 \(\rho_{\rm obs}=(1-v)\rho_0\)，故
\[
v=1-\frac{10.20}{10.26482}=0.00631484\approx0.6315\%.
\]
计算只在“体密度差异仅由空位造成、晶胞常数保持不变”的指定模型中成立，不应将此数值当作真实银晶体的实验空位浓度。
\begin{reviewnote}
本题计算机独立数值见证：1：0.23771491；2：0.41173430；3：10.264821；4：0.0063148443。
逆向方程最大相对残差：0.000e+0。
\end{reviewnote}
\chapter{分析化学：解答}
\section{A01 \quad 双波长光谱解混、稀释反算与试样质量}
由吸光度加和关系，
\[
\begin{bmatrix}10000&2000\\3000&8000\end{bmatrix}
\begin{bmatrix}c_X\\c_Y\end{bmatrix}
=\begin{bmatrix}0.4600\\0.3400\end{bmatrix}.
\]
行列式为 \(D=10000\times8000-3000\times2000=74.0\times10^6\neq0\)，故
\[
c_X=\frac{8000(0.4600)-2000(0.3400)}D
    =4.054054\times10^{-5}\,\mathrm M,
\]
\[
c_Y=\frac{10000(0.3400)-3000(0.4600)}D
    =2.729730\times10^{-5}\,\mathrm M.
\]
定容导致 \(10.00\) 倍稀释，原样浓度分别为 \(4.054054\times10^{-4}\) 与 \(2.729730\times10^{-4}\,\mathrm M\)。取 \(25.00\,\mathrm{mL}\) 原样，
\[
m_X=c_{X,0}(0.02500)(150.0)=1.52027\times10^{-3}\,\mathrm g,
\quad
m_Y=c_{Y,0}(0.02500)(200.0)=1.36486\times10^{-3}\,\mathrm g.
\]
代回两波长方程应分别恢复 \(0.4600\) 与 \(0.3400\)。
\begin{reviewnote}
本题计算机独立数值见证：1：0.000040540541；2：0.000027297297；3：0.00040540541；4：0.00027297297；5：0.0015202703；6：0.0013648649。
逆向方程最大相对残差：1.207e-16。
\end{reviewnote}
\section{A02 \quad 双指示剂滴定反演碳酸盐组成与离子强度}
设取样中的 \(\ce{CO3^{2-}}\) 与 \(\ce{HCO3-}\) 物质的量为 \(x,y\)。第一终点仅有
\(\ce{CO3^{2-} -> HCO3-}\)，故 \(x=0.1000(0.00820)=0.0008200\,\mathrm{mol}\)。
第二终点累计耗酸满足 \(2x+y=0.1000(0.02160)\)，得到
\(y=0.0005200\,\mathrm{mol}\)。因此
\[
c_{\ce{Na2CO3}}=\frac{x}{0.02500}=0.03280\,\mathrm M,\qquad
c_{\ce{NaHCO3}}=\frac{y}{0.02500}=0.02080\,\mathrm M.
\]
以 \(M(\ce{Na2CO3})=105.99\)、\(M(\ce{NaHCO3})=84.0066\,\mathrm{g\,mol^{-1}}\)，
\[
m_1=0.03280(0.2500)(105.99)=0.86912\,\mathrm g,\quad
m_2=0.02080(0.2500)(84.0066)=0.43683\,\mathrm g.
\]
形式电荷平衡给 \([\ce{Na+}]=2c_1+c_2\)，所以完全解离模型下
\[
I=\tfrac12\bigl([\ce{Na+}]+4c_1+c_2\bigr)=3c_1+c_2
=0.11920\,\mathrm M.
\]
该 \(I\) 是按给定形式浓度的估算；若考虑酸碱物种与活度，需重解平衡。
\begin{reviewnote}
本题计算机独立数值见证：1：0.032800000；2：0.020800000；3：0.86911800；4：0.43683432；5：0.11920000。
逆向方程最大相对残差：0.000e+0。
\end{reviewnote}
\section{A03 \quad 钙镁络合差量滴定、空白校正与总硬度}
空白校正体积为 \(V_{\rm Ca}=17.40\,\mathrm{mL}\)、\(V_{\rm tot}=28.20\,\mathrm{mL}\)。由 EDTA 的 \(1:1\) 络合计量，
\[
c_{\rm Ca}=\frac{0.01000(0.01740)}{0.05000}
=0.003480\,\mathrm M,
\]
\[
c_{\rm Mg}=\frac{0.01000(0.02820-0.01740)}{0.05000}
=0.002160\,\mathrm M.
\]
若有同浓度 \(0.5000\,\mathrm L\) 样品，则其中金属质量为
\[
m_{\rm Ca}=c_{\rm Ca}(0.5000)(40.078)=0.069736\,\mathrm g,\qquad
m_{\rm Mg}=c_{\rm Mg}(0.5000)(24.305)=0.0262494\,\mathrm g.
\]
将两者的摩尔浓度和以相同 \(2+\) 价的 \(\ce{CaCO3}\) 当量换算，
\[
H=(c_{\rm Ca}+c_{\rm Mg})(100.09)(1000)
=564.51\,\mathrm{mg\,L^{-1}}.
\]
这里的硬度是 \emph{以 \(\ce{CaCO3}\) 计}，不是实际生成的碳酸钙质量。
\begin{reviewnote}
本题计算机独立数值见证：1：0.0034800000；2：0.0021600000；3：0.069735720；4：0.026249400；5：564.50760。
逆向方程最大相对残差：2.042e-16。
\end{reviewnote}
\chapter{有机化学：解答}
\section{O01 \quad 连续一级反应的中间体最佳分离时刻}
设 \(a,b,c\) 为对初始 \(A\) 物质的量归一化的分数。动力学方程
\[
a'=-k_1a,\qquad b'=k_1a-k_2b,\qquad c'=k_2b,\qquad
(a,b,c)(0)=(1,0,0).
\]
线性微分方程给出
\[
a(t)=e^{-k_1t},\qquad
b(t)=\frac{k_1}{k_2-k_1}
\left(e^{-k_1t}-e^{-k_2t}\right),\qquad
c(t)=1-a(t)-b(t).
\]
求导 \(b'(t)=0\)，等价于 \(k_2e^{-k_2t}=k_1e^{-k_1t}\)，故
\[
t_*=\frac{\ln(k_2/k_1)}{k_2-k_1}
=15.2715\,\mathrm h.
\]
该时刻
\[
(a,b,c)=(0.217153,\ 0.542884,\ 0.239963).
\]
因此 \(B\) 的可分离理论质量（若瞬时完美分离）为
\[
m_B=(0.2500\,\mathrm{mol})\times0.542884\times
150.0\,\mathrm{g\,mol^{-1}}=20.3581\,\mathrm g.
\]
该极值唯一，因为 \(b(0)=0\)、\(b(t)>0\) 且 \(b(\infty)=0\)。
\begin{reviewnote}
本题计算机独立数值见证：1：15.271512；2：0.21715341；3：0.54288352；4：0.23996307；5：20.358132。
逆向方程最大相对残差：0.000e+0。
\end{reviewnote}
\section{O02 \quad 非等当量缩聚的 Carothers 关系、脱水守恒与数均分子量}
双官能起始单体的分子 mol 数为
\[
N_0=\frac{0.9600+1.000}{2}=0.9800\,\mathrm{mol}.
\]
每形成一条键消耗一个 A 与一个 B 官能团，故反应键数
\[
\xi=pN_{A,0}=0.9700(0.9600)=0.9312\,\mathrm{mol}.
\]
无环化时每形成一条键，独立分子数减少一，最终链数
\(N_{\rm chain}=0.9800-0.9312=0.04880\,\mathrm{mol}\)。因而
\[
X_n=\frac{N_0}{N_{\rm chain}}
=\frac{1+r}{1+r-2rp}
=20.08197,\qquad r=\frac{N_{A,0}}{N_{B,0}}=0.9600.
\]
初始总质量 \(m_0=(0.9800)(130.0)=127.4\,\mathrm g\)，水损失 \(m_{\rm water}=0.9312(18.015)=16.775568\,\mathrm g\)，
\[
m_{\rm polymer}=110.624432\,\mathrm g,\qquad
M_n=\frac{m_{\rm polymer}}{N_{\rm chain}}
=2266.89\,\mathrm{g\,mol^{-1}}.
\]
不能简单写成 \(X_n\times130.0\)：因为每条缩合键有实际脱水质量损失。
\begin{reviewnote}
本题计算机独立数值见证：1：0.98000000；2：0.93120000；3：0.048800000；4：20.081967；5：110.62443；6：2266.8941。
逆向方程最大相对残差：1.285e-16。
\end{reviewnote}
\section{O03 \quad 双响应因子手性色谱的 ee、质量平衡与收率}
色谱面积不能直接用于求质量比，须先除以各自的响应因子：
\[
m_R=\frac{1500}{1.20}=1250\,\mathrm{mg},\qquad
m_S=\frac{700}{0.800}=875\,\mathrm{mg}.
\]
由于两对映体摩尔质量相同，物质的量比例等于质量比例。总分离产物 \(m=2125\,\mathrm{mg}\)，\(R\) 为主对映体，其质量百分数为
\[
w_R=100\frac{1250}{2125}=58.8235\%.
\]
对映体过量和总收率分别为
\[
ee=100\frac{1250-875}{1250+875}=17.6471\%,\qquad
Y=100\frac{2125}{2500}=85.00\%.
\]
“ee”衡量立体选择性；“总收率”衡量实际回收量相对理论产量。二者不可混用。
\begin{reviewnote}
本题计算机独立数值见证：1：1250.0000；2：875.00000；3：58.823529；4：17.647059；5：85.000000。
逆向方程最大相对残差：0.000e+0。
\end{reviewnote}
\chapter{物理化学：解答}
\section{P01 \quad 双层扩散膜的等效阻力、通量与累积转移质量}
各层稳态 Fick 定律 \(J=-D_i\,dc_i/dx\)；按膜--水分配系数把两层的浓度降换算到同一水相参照浓度，可得到串联阻力
\[
\mathcal R=\frac{L_1}{D_1K_1}+\frac{L_2}{D_2K_2}
=\frac{80\times10^{-6}}{1.00\times10^{-10}(0.500)}
+\frac{120\times10^{-6}}{3.00\times10^{-10}(0.800)}
=2.10\times10^6\,\mathrm{s\,m^{-1}}.
\]
因此
\[
J=\frac{c_{\mathrm L}-c_{\mathrm R}}{\mathcal R}
=\frac{1.800}{2.10\times10^6}
=8.57143\times10^{-7}\,\mathrm{mol\,m^{-2}s^{-1}}.
\]
面积 \(A=2.00\times10^{-4}\,\mathrm{m^2}\)、时间 \(t=7200\,\mathrm s\)，故
\[
n=JAt=1.23429\times10^{-6}\,\mathrm{mol},\qquad
m=n(180.0)\times1000=0.222171\,\mathrm{mg}.
\]
保持两储液浓度不变是恒定通量计算的必要外部条件。
\begin{reviewnote}
本题计算机独立数值见证：1：2100000.0；2：8.5714286e-7；3：0.0000012342857；4：0.22217143。
逆向方程最大相对残差：0.000e+0。
\end{reviewnote}
\section{P02 \quad 热容差修正反应热力学与平衡常数}
由 Kirchhoff 定律和熵的温度积分，
\[
\Delta_rH^\circ(T)=H_0+\Delta C_p(T-T_0),\qquad
\Delta_rS^\circ(T)=S_0+\Delta C_p\ln(T/T_0).
\]
取 \(T=350.00\,\mathrm K\)，得到
\[
\Delta_rH^\circ=47.5925\,\mathrm{kJ\,mol^{-1}},
\quad
\Delta_rS^\circ=128.01682\,\mathrm{J\,mol^{-1}K^{-1}}.
\]
应先把焓与 \(T\Delta S\) 统一成 J/mol，计算
\[
\Delta_rG^\circ=47592.5-350(128.01682)
=2786.61\,\mathrm{J\,mol^{-1}}.
\]
按标准活度定义的无量纲 \(K\) 满足
\[
K=\exp\!\left(-\frac{\Delta_rG^\circ}{RT}\right)
=0.383821,\qquad
\left.\frac{d\ln K}{dT}\right|_{350}
=\frac{\Delta_rH^\circ(350)}{R(350)^2}
=0.0467270\,\mathrm{K^{-1}}.
\]
最后一式可由前述 \(H(T),S(T)\) 的显式表达求导复核。
\begin{reviewnote}
本题计算机独立数值见证：1：47592.500；2：128.01682；3：2786.6123；4：0.38382124；5：0.046727037。
逆向方程最大相对残差：0.000e+0。
\end{reviewnote}
\section{P03 \quad 有限溶液 Langmuir 吸附的非线性质量平衡}
平衡时总物料守恒为
\[
VC_0=VC+mq_{\max}\frac{KC}{1+KC}.
\]
化为二次多项式，
\[
VK C^2+\bigl[V+m q_{\max}K-VC_0K\bigr]C-VC_0=0.
\]
因常数项为负，二次方程恰有一个正根；等价地，原始右侧对 \(C\ge0\) 严格递增，故物理解唯一。数值代入并取正根，
\[
C=0.001354066\,\mathrm{mol\,L^{-1}}.
\]
相应表面覆盖度
\[
\theta=\frac{KC}{1+KC}=0.4037088,
\]
吸附总物质的量
\[
n_{\rm ads}=m q_{\max}\theta
=3.22967\times10^{-4}\,\mathrm{mol},
\]
去除率
\[
\eta=100\left(1-\frac{C}{C_0}\right)=32.2967\%.
\]
可代回验证 \(VC+n_{\rm ads}=VC_0=0.001000\,\mathrm{mol}\)，不允许用初始浓度直接计算最终 \(\theta\)。
\begin{reviewnote}
本题计算机独立数值见证：1：0.0013540659；2：0.40370880；3：0.00032296704；4：32.296704。
逆向方程最大相对残差：2.168e-16。
\end{reviewnote}
\chapter*{质量控制与发布边界}
\addcontentsline{toc}{chapter}{质量控制与发布边界}
\begin{enumerate}
\item 12 道题各有独立内容，不使用批量参数变式增加“独立题型”计数。
\item 每题列出前提、关键守恒关系和数值反算；计算过程可以由版本固定的脚本重现。
\item 数值反算无法自动检验文献真实性、实验可行性、结构化学细节和教学适切性。
\item 所有题目仍为 awaiting\_expert\_review、publishable=false，
不自动进入正式组卷索引。
\item PDF 为专家审校稿。若修改题意、反应模型或常数，必须同时重新生成 PDF、
题目 JSON 与科学审核报告。
\end{enumerate}
\end{document}
